\documentclass[11pt,a4paper]{article}
\usepackage[utf8]{inputenc}
\usepackage[T1]{fontenc}
\usepackage[brazilian]{babel}
\usepackage{geometry}
\usepackage{graphicx}
\usepackage{booktabs}
\usepackage{float}
\usepackage{enumitem}
\usepackage{amsmath}
\usepackage{xcolor}
\usepackage{tcolorbox}
\geometry{margin=2.3cm}
\setlength{\parindent}{0pt}
\setlength{\parskip}{0.5em}
\tcbuselibrary{skins,breakable}
\newtcolorbox{nota}[1][]{colback=yellow!10,colframe=orange!60!black,
  fonttitle=\bfseries,title=#1,breakable}

\begin{document}
\begin{center}
  
\includegraphics[width=0.16\textwidth]{../../../figures_global/unb.png}\\[0.3cm]
  {\large\textbf{GABARITO DO PROFESSOR -- USO RESTRITO}}\\[0.2cm]
  {\large FGA0092 -- Princípios de Comunicação para Engenharia}\\[0.2cm]
  \rule{12cm}{0.8pt}\\[0.2cm]
  {\Large\textbf{Prática 01: FFT e Análise Espectral}}\\[0.2cm]
  \rule{12cm}{0.8pt}
\end{center}

\begin{nota}[Origem dos valores]
Todos os números deste gabarito foram gerados por
\texttt{Codigos/pratica\_01\_gabarito/codigos/pratica\_01\_gabarito.py}, que
executa o mesmo \textit{flowgraph} da prática em modo headless e calcula o
espectro do mesmo modo que o \texttt{QT GUI Frequency Sink}, ou seja
$X_{\mathrm{dB}}[k] = 20\log_{10}(|X[k]|/N)$, com as janelas do próprio GNU
Radio e média de potência entre blocos, equivalente a
\textit{Average: High}. Os CSVs correspondentes estão em
\texttt{Codigos/pratica\_01\_gabarito/dados/}.

Aceitar tolerância de $\pm 2$\,dB nas leituras feitas na tela e de $\pm 1$
índice espectral nas leituras de frequência. O eixo do bloco é
\textit{Relative Gain}, de referência arbitrária, portanto avaliar apenas
diferenças de níveis, nunca valores absolutos.
\end{nota}

% ==========================================================
\section*{Gráficos esperados}

As figuras abaixo são geradas por \texttt{pratica\_01\_gabarito.py} e ficam em
\texttt{Codigos/pratica\_01\_gabarito/figuras/}. Servem de referência para
comparar com as capturas de tela entregues pelo aluno.

\begin{figure}[H]\centering
  \includegraphics[width=0.92\textwidth]{../../Codigos/pratica_01_gabarito/figuras/cenario_C_dois_tons_ruido.png}
  \caption*{Cenário C do enunciado: dois tons com \texttt{noise\_amp} $= 5{,}0$.
    O tom de 2500\,Hz já não satisfaz o critério de visibilidade.}
\end{figure}

\begin{figure}[H]\centering
  \includegraphics[width=0.92\textwidth]{../../Codigos/pratica_01_gabarito/figuras/etapa6_janelas.png}
  \caption*{Etapa 6 com o tom em 1050\,Hz. A diferença entre a janela retangular
    e a Blackman-harris passa de 70\,dB nas saias de vazamento. Com o tom em
    1000\,Hz os quatro painéis ficariam indistinguíveis.}
\end{figure}

\clearpage

% ==========================================================
\section*{Tabelas de referência}

\subsection*{Etapa 2 -- posição do pico}

\begin{table}[H]\centering
\begin{tabular}{c c c c}
  \toprule
  $f_1$ (Hz) & $k = f_1 N/f_s$ & $f_k$ (Hz) & Pico lido (Hz) \\
  \midrule
  500  & 16{,}0  & 500{,}0  & 500 \\
  1000 & 32{,}0  & 1000{,}0 & 1000 \\
  2000 & 64{,}0  & 2000{,}0 & 2000 \\
  4000 & 128{,}0 & 4000{,}0 & 4000 \\
  \bottomrule
\end{tabular}
\caption*{As quatro frequências são múltiplos exatos de $\Delta f = 31{,}25$\,Hz,
logo $k$ é inteiro e o pico lido coincide com $f_1$.}
\end{table}

\subsection*{Etapa 4 -- resolução espectral, tons de 1000 e 1200\,Hz}

\begin{table}[H]\centering
\setlength{\tabcolsep}{5pt}
\begin{tabular}{c c c c c c c}
  \toprule
  $N$ & $\Delta f$ (Hz) & $\Delta k$ & $A_{\min,\mathrm{dB}}$ &
    $V_{\mathrm{dB}}$ & $D_{\mathrm{dB}}$ & $D_{\mathrm{dB}} \geq 3$\,dB? \\
  \midrule
  256  & 125{,}00 & 1{,}60  & $-12{,}6$ & $-13{,}0$  & 0{,}3  & Não \\
  512  & 62{,}50  & 3{,}20  & $-12{,}3$ & $-21{,}0$  & 8{,}7  & Sim \\
  1024 & 31{,}25  & 6{,}40  & $-13{,}0$ & $-60{,}6$  & 47{,}6 & Sim \\
  2048 & 15{,}62  & 12{,}80 & $-12{,}3$ & $-85{,}4$  & 73{,}2 & Sim \\
  4096 & 7{,}81   & 25{,}60 & $-12{,}9$ & $-101{,}3$ & 88{,}4 & Sim \\
  \bottomrule
\end{tabular}
\caption*{A transição ocorre entre $N = 256$ e $N = 512$. Para $N \geq 1024$ o
fundo do vale desce muito e o valor lido varia bastante com a média
selecionada; aceitar qualquer valor que caracterize $D_{\mathrm{dB}} \gg 3$\,dB.}
\end{table}

\subsection*{Etapa 5 -- efeito do ruído}

\begin{table}[H]\centering
\setlength{\tabcolsep}{5pt}
\begin{tabular}{c c c c c c}
  \toprule
  $\sigma$ & $\mathrm{SNR}_{1,\mathrm{dB}}$ & $\mathrm{SNR}_{2,\mathrm{dB}}$ &
    $P_{\mathrm{dB}}$ & $A_{2,\mathrm{dB}} - P_{\mathrm{dB}}$ & Tom 2 visível? \\
  \midrule
  0{,}5  & 3{,}0     & $-3{,}0$  & $-40{,}4$ & 22{,}3 & Sim \\
  1{,}0  & $-3{,}0$  & $-9{,}0$  & $-34{,}4$ & 16{,}4 & Sim \\
  2{,}0  & $-9{,}0$  & $-15{,}1$ & $-28{,}3$ & 10{,}6 & Sim \\
  5{,}0  & $-17{,}0$ & $-23{,}0$ & $-20{,}4$ & 4{,}2  & Não \\
  10{,}0 & $-23{,}0$ & $-29{,}0$ & $-14{,}4$ & 1{,}4  & Não \\
  \bottomrule
\end{tabular}
\caption*{Critério de visibilidade: $\mathrm{SNR}_{\mathrm{disp,dB}} \geq 6$\,dB.
O tom 2 cruza o critério em $\sigma \approx 3{,}3$.}
\end{table}

\subsection*{Etapa 6 -- funções de janela, tom em 1050\,Hz}

\begin{table}[H]\centering
\begin{tabular}{l c c c c}
  \toprule
  Janela & $A_{\mathrm{pico,dB}}$ & $L_{\mathrm{dB}}$ & $S_{\mathrm{dB}}$ &
    Largura a $-3$\,dB (Hz) \\
  \midrule
  Rectangular     & $-8{,}4$  & $-36{,}7$  & 28{,}2 & 35{,}2 \\
  Hamming         & $-12{,}5$ & $-59{,}1$  & 46{,}6 & 40{,}1 \\
  Hann            & $-13{,}0$ & $-83{,}4$  & 70{,}5 & 42{,}6 \\
  Blackman-harris & $-15{,}5$ & $-114{,}5$ & 99{,}0 & 56{,}2 \\
  \bottomrule
\end{tabular}
\caption*{As duas últimas colunas se ordenam em sentidos opostos, que é o
resultado central da etapa. $L_{\mathrm{dB}}$ é lido fora de
$f_1 \pm 300$\,Hz.}
\end{table}

% ==========================================================
\section*{Respostas esperadas}

\begin{enumerate}[label=\textbf{Q\arabic*.}, leftmargin=*, itemsep=1em]

  \item \textbf{Posição do pico e a relação $f_k = k f_s/N$.}

    Espera-se o cálculo explícito $k = f_1 N / f_s$. Para $f_1 = 2000$\,Hz,
    $N = 1024$ e $f_s = 32000$\,Hz resulta $k = 64$, e
    $f_k = 64 \cdot 32000/1024 = 2000$\,Hz. O aluno deve reconhecer que a
    transformada só fornece amostras do espectro em múltiplos de
    $\Delta f = 31{,}25$\,Hz. Uma frequência que não seja múltipla de
    $\Delta f$ é atribuída ao índice mais próximo, o que limita o erro de
    leitura a $\Delta f/2 = 15{,}6$\,Hz.

    \begin{nota}[Erros comuns]
      Tratar bins e hertz como a mesma grandeza sem aplicar a conversão.
      Afirmar que a resolução depende apenas de $N$, esquecendo $f_s$. Não
      justificar o limite $\Delta f/2$.
    \end{nota}

  \item \textbf{Menor $N$ pelo critério $D_{\mathrm{dB}} \geq 3$\,dB e extensão
    para 100\,Hz.}

    Da tabela da Etapa 4, o critério passa a ser satisfeito em $N = 512$, com
    $\Delta k = 3{,}20$. Em $N = 256$ tem-se $\Delta k = 1{,}60$: os dois tons
    caem dentro do mesmo lóbulo principal da janela Hann, que ocupa cerca de 4
    índices, e por isso $D_{\mathrm{dB}} = 0{,}3$\,dB.

    Exigindo o mesmo $\Delta k = 3{,}20$ para tons separados de 100\,Hz,
    \[
      N \;\geq\; \Delta k \cdot \frac{f_s}{100}
        \;=\; 3{,}20 \cdot \frac{32000}{100} \;=\; 1024 ,
    \]
    ou seja $N = 1024$. O limite ingênuo $\Delta f \leq 100$\,Hz daria
    $N \geq 320$, portanto $N = 512$, que é insuficiente: ele garante mais de
    um índice de separação, mas não a profundidade de vale exigida, porque
    ignora a largura do lóbulo principal da janela.

    \begin{nota}[Erros comuns]
      Responder $N = 512$ usando apenas $\Delta f \leq 100$\,Hz. Essa resposta
      vale no máximo metade da pontuação: o critério do enunciado é a
      profundidade do vale, não o espaçamento de índices. Também é erro
      preencher a coluna $D_{\mathrm{dB}}$ sem apresentar as leituras de
      $A_{\min,\mathrm{dB}}$ e $V_{\mathrm{dB}}$ que a originaram.
    \end{nota}

  \item \textbf{Superposição dos dois tons.}

    A diferença medida entre os picos deve ser
    $20\log_{10}(\mathrm{amp1}/\mathrm{amp2}) = 20\log_{10}(2) = 6{,}0$\,dB.
    A justificativa é a linearidade da transformada: $X[k] = X_1[k] + X_2[k]$.
    Como as duas frequências caem em índices distintos, cada pico conserva o
    nível que teria isoladamente. Com um elemento não linear no caminho
    apareceriam harmônicos e produtos de intermodulação, por exemplo em
    1500\,Hz e 3500\,Hz, que não são observados aqui.

    \begin{nota}[Erros comuns]
      Afirmar que os espectros se combinam por convolução: a convolução ocorre
      no produto de sinais, não na soma. Dizer apenas que ``aparecem dois
      picos'' sem verificar a diferença de 6\,dB.
    \end{nota}

  \item \textbf{Limiar de visibilidade e ganho de processamento.}

    Da tabela da Etapa 5, o tom 2 satisfaz o critério até $\sigma = 2{,}0$ e o
    perde em $\sigma = 5{,}0$. Como $A_{2,\mathrm{dB}} - P_{\mathrm{dB}}$ cai
    $20\log_{10}\sigma$, o limiar exato é $\sigma \approx 3{,}3$.

    Em $\sigma = 1{,}0$ a SNR de banda larga do tom 2 vale $-9{,}0$\,dB,
    enquanto a SNR do display vale $16{,}4$\,dB. A diferença de $25{,}4$\,dB é o
    ganho de processamento da FFT: o tom concentra sua energia em um único
    índice, enquanto o ruído se reparte pelos $N/2$ índices exibidos. Dobrar o
    \textit{FFT Size} baixa o piso em 3\,dB e eleva a SNR do display em 3\,dB,
    sem alterar a SNR de banda larga. É por isso que um tom com SNR de banda
    larga negativa ainda aparece com clareza na tela.

    \begin{nota}[Erros comuns]
      Confundir \texttt{noise\_amp} com SNR, ou tratá-lo como porcentagem. Foi
      o erro mais frequente em 2026/1. Responder apenas ``o piso sobe'' sem
      nenhum valor em dB. Não perceber que as duas SNR são grandezas
      diferentes e que a do display depende de $N$.
    \end{nota}

  \item \textbf{Compromisso entre supressão e largura.}

    Ordem por supressão, do maior para o menor:
    Blackman-harris (99\,dB), Hann (70\,dB), Hamming (47\,dB),
    Rectangular (28\,dB). Ordem por largura do lóbulo principal, da menor para
    a maior: Rectangular (35\,Hz), Hamming (40\,Hz), Hann (43\,Hz),
    Blackman-harris (56\,Hz). As duas ordens são inversas.

    A janela retangular preserva o maior nível de pico, $-8{,}4$\,dB, e a
    Blackman-harris oferece a maior supressão. As duas propriedades não podem
    ser maximizadas ao mesmo tempo porque suprimir lóbulos laterais exige
    suavizar as bordas do bloco, o que reduz $\sum_n w[n]$ e alarga o lóbulo
    principal.

    A Etapa 6 usa 1050\,Hz porque 1000\,Hz cai exatamente sobre o índice
    $k = 32$: nesse caso o bloco de $N$ amostras contém um número inteiro de
    períodos, não há descontinuidade nas bordas e o vazamento desaparece no
    piso do gráfico, tornando as quatro janelas indistinguíveis. Em 1050\,Hz o
    índice seria $k = 33{,}6$ e o vazamento fica cerca de 28\,dB abaixo do pico
    com janela retangular.

    \begin{nota}[Erros comuns]
      Dizer que a retangular é ``sempre pior''. Comparar janelas sem valores em
      dB. Não responder por que a etapa usa 1050\,Hz, que é o ponto conceitual
      da questão.
    \end{nota}

\end{enumerate}

% ==========================================================
\section*{Critérios de correção}

\begin{table}[H]\centering
\begin{tabular}{l c}
  \toprule
  Critério & Pontuação \\
  \midrule
  Fundamentação teórica com as equações da prática        & 1{,}0 \\
  Tabelas da Etapa 2 e da Etapa 4 preenchidas             & 1{,}0 \\
  Tabelas da Etapa 5 e da Etapa 6 preenchidas             & 1{,}0 \\
  Capturas dos três cenários, no tempo e na frequência    & 1{,}0 \\
  Q1 -- cálculo de $k$ e limite $\Delta f/2$              & 1{,}0 \\
  Q2 -- menor $N$ pelo critério do vale e extensão        & 1{,}0 \\
  Q3 -- linearidade e diferença de 6\,dB verificada       & 1{,}0 \\
  Q4 -- limiar de $\sigma$ e ganho de processamento       & 1{,}0 \\
  Q5 -- compromisso entre supressão e largura             & 1{,}0 \\
  Conclusão e organização do relatório                    & 1{,}0 \\
  \midrule
  \textbf{Total} & \textbf{10{,}0} \\
  \bottomrule
\end{tabular}
\end{table}

\begin{nota}[Nota de correção geral]
Priorizar a evidência experimental. Resposta sem o valor medido que a sustenta
não recebe pontuação integral, mesmo que o raciocínio esteja correto. Não
penalizar diferenças de até $\pm 2$\,dB em relação a este gabarito, nem o uso
de uma média diferente no painel de controle.

Em 2026/1 a fundamentação teórica teve 12 zeros em 29 relatórios, porque o
modelo entregue ao aluno não tinha essa seção. Confirmar que o modelo divulgado
no SIGAA contém uma seção de fundamentação teórica antes de aplicar esse
critério.
\end{nota}

\end{document}
